\section{Findings}
\label{sec:findings}

Stating the laws and then checking them found six defects (1--6), and a
subsequent audit of every place that compares names or paths, and of the
environment assumptions behind the rules, found ten more (7--16).  A second
independent review, of the Quick Look endpoint added later, found two more (17--18)
against Laws~\ref{law:absent} and~\ref{law:indist}.  Each was
reproduced by a failing test before it was fixed.  Findings 1, 7, 8, 10 to 15, 17
and 18 are security defects; the rest are consistency defects that fail closed
or are cosmetic, but each is a place where two parts of the program disagreed.

\begin{longtable}{@{}p{0.03\linewidth}p{0.28\linewidth}p{0.23\linewidth}p{0.35\linewidth}@{}}
\toprule
& Defect & Law broken & Fix \\
\midrule
\endhead
1 & \textbf{Deny bypass through full case folding.}  On APFS the name
\code{.\ss h} opens \code{.ssh}, but the matcher folded with simple folding, so
a rule for \code{\~{}/.ssh} did not match it; \code{Open}, \code{Head} and the
listing of the parent all served the denied folder.  Also long s, the Kelvin
sign and the ligatures.
& \ref{law:fold-invariance}, \ref{law:absent}
& Rules, roots and search compare after full folding (\code{rules.Fold}); the alias
set was measured on this file system and is a test.  \label{find:fold} \\
2 & \textbf{Listed but unopenable.}  A plain file named like a \code{dir/} deny
rule (a file called \code{cache}) was shown by \code{List} but refused by
\code{Open}, because \code{Open}'s pre-check assumed every path is a directory.
& \ref{law:listed-open}
& The pre-check treats the final component as a file; the precise check after
opening still refuses real directories. \\
3 & \textbf{Glob class matched across a slash.}  \code{[!a]} compiled to
\code{[\^{}a]}, which matches \code{/}, so a pattern with no slash could match
across directory components.
& \ref{law:component}
& Negated classes exclude \code{/}; a \code{/} inside a class is a pattern error. \\
4 & \textbf{Row order not transitive.}  A modified time that could not be parsed
compared equal to everything, so \code{a}$=$\code{b}, \code{a}$=$\code{c} but
\code{b}$<$\code{c}.  The resulting order depended on the input order.
& \ref{law:total}, \ref{law:input-order}
& Unreadable times sort as the oldest; comparisons no longer use subtraction. \\
5 & \textbf{Filter and search disagreed.}  The filter box folded with simple
lower-casing, so \code{strasse} did not find \code{Stra\ss e}, while filename
search did.
& \ref{law:filter}
& The filter folds with lower-, upper-, lower-casing, which matches full folding
on the measured aliases. \\
6 & \textbf{Stale fallback.}  After an image failed to load, the fall-back request
had its own cancellation handle, so a slow answer could replace the preview of a
newer selection.  \label{find:stale}
& \ref{law:fresh}
& One handle per selection is shared by every request started for it. \\
7 & \textbf{Firmlink prefix matched by byte length.}  The data-volume prefix and
the firmlink directories were recognized by slicing the path to the prefix's
length, so a folded spelling (\code{/S}long-s\code{ystem}, or \code{Data} with
the \code{st} ligature) was cut mid-character and not reduced.  Under
\code{--root /} the unreduced path was inside the root and matched no rule.
& \ref{law:alias-fold}, \ref{law:absent}
& Components are compared one at a time with \code{Fold}. \\
8 & \textbf{Root containment by spelling.}  Containment folded names, so on a
case-sensitive volume a root \code{Public} also served the different folder
\code{public}.  Folding is safe for deny rules but not for an allow check.
& \ref{law:root-identity}, \ref{law:containment}
& Containment is decided by identity: the corresponding leading directory must be
the root itself. \\
9 & \textbf{Client and CLI comparisons.}  Breadcrumbs sliced the path by the
root's character count (wrong when folding changes length) and compared roots
case-sensitively; a link or hash naming a file with another case or accent
selected nothing; the core-rule warning compared rule text exactly.
& \ref{law:roots-prefix}
& All use the same name equality; breadcrumbs split by components. \\
10 & \textbf{Credential rules anchored to one home directory.}  The rules for
\code{\~{}/.ssh}, \code{\~{}/.aws} and the others named the current home only, so
the same folders in a backup or clone of the home folder, in another account's
home, or under a wrong \code{\$HOME} were served.
& \ref{law:credential-anywhere}
& The credential rules match at any depth, like the secret-file patterns. \\
11 & \textbf{A path is not a file (hard links).}  Rules name paths, so a hard link
to a credential file placed anywhere else was served under its innocent name; this
had been recorded only as a limitation.
& \ref{law:absent}, \ref{law:linked}
& A file with several names is judged by identity: a denied name denies them all, in \code{Open} and in listings.  First done for the credential locations only. \\
12 & \textbf{The session cookie followed the host, not the port.}  Cookies are
scoped by host, so the cookie was also sent to every other web server on
\code{127.0.0.1} that the same browser visited, which could replay it.  It had
been recorded as a limitation.
& \ref{law:cookie-path}
& Everything is served under a random per-launch prefix and the cookie's path is
that prefix; requests outside it are refused. \\
13 & \textbf{Hard links to \code{.env}, \code{.pem} and \code{.key} files.}  The
identity check of finding 11 covered only the credential folders, so a second name
for a secret file elsewhere (\code{notes.txt} linked to a \code{.env}) was still
served.
& \ref{law:linked}
& The check covers every denied name under the roots and protected locations, using
a background index of multi-name files; multi-name files are refused until the
index is ready. \\
14 & \textbf{Rule files that load and do nothing.}  A byte order mark broke the
first rule; bare carriage returns joined lines; an inline comment, \code{\~{}user/},
quotes, an invisible character or a lookalike slash produced a rule that never
matched.  All loaded without an error.
& \ref{law:rule-file}
& The byte order mark and line endings are normalized; the rest are errors that
say what is wrong, so fsb does not start. \\
15 & \textbf{Secrets under other names.}  Only known folders and the extensions
\code{.pem} and \code{.key} were denied, so a private key copied out of
\code{.ssh} (\code{id\_rsa}), a certificate store or a password database was served.
& \ref{law:credential-anywhere}
& The names that identify such files are core rules at any depth (\code{.pub} files
stay reachable). \\
16 & \textbf{A name that shows as another name.}  A right-to-left override or a
control character in a file name reordered or hid what the listing displayed, so
\code{invoice}\textit{U+202E}\code{txt.exe} read as a text file.
& \ref{law:display}
& Deceptive characters are replaced, for display only, by a visible marker; real
right-to-left text is untouched and downloads keep the real name. \\
17 & \textbf{A path that means more to another program.}  Quick Look was given the
real path of a document that passed every check, but a Finder alias named
\code{report.docx} is a regular file that Quick Look resolves to its target, so it
drew a file in a denied folder, or outside the roots.
& \ref{law:absent}
& Quick Look is given only a private copy of bytes read through $\Open$
(Law~\ref{law:quicklook}); a real alias is the regression test. \\
18 & \textbf{A package answered differently for a hidden part.}  A package document
holding a denied file was refused (404) while a clean one answered 200, although
listings show both alike.
& \ref{law:indist}
& Parts a listing leaves out are left out of the copy, and the package answers as a
clean one. \\
\bottomrule
\end{longtable}

\subsection{Observations that are not defects}

\begin{itemize}
  \item A dangling symbolic link is listed (marked as broken) but cannot be
        opened.  Law~\ref{law:listed-open} excludes it.
  \item Hidden entries are not listed but can be opened, and a hidden directory
        that is entered on purpose shows its contents.
        (Law~\ref{law:hidden-reachable}.)
  \item A symbolic link \emph{named} like a denied file is refused even when its
        target is harmless.  Definition~\ref{def:ok} checks both names.
  \item The macOS by-inode path \code{/.vol/<device>/<inode>} was probed as a
        possible alias that would evade path rules; on the tested system the
        directory is empty and the path cannot be opened.
\end{itemize}

\subsection{Defects found afterwards by review}

Reviews of the code (October 2026) found defects that these laws did not, each in
a place the laws state only loosely: a hard-link verdict served for a file whose
names the index had not seen, or whose names changed (Law~\ref{law:linked} now says
what ``determined'' means); a search that entered a folder swapped for a link after it
was queued, and that looked complete when a folder could not be opened (the search
definition now says so); a collator that left two names in input order (the row order
now ends on the raw characters); and several limits that did not bind (a count that could
never trip, a work limit that ignored excluded entries).  The design specification's
history lists them; each has a test that fails without its fix.  Later reviews of
opening the browser on Linux found that the launch URL, in the opener's command line,
let another user take the session; Law~\ref{law:launch-owner} was then stated, and its
review found that a used token still caused a lookup and a warning on every replay (the
law's first case now comes before the owner), and that the tests could not tell the
client's end of a connection from \texttt{fsb}'s own, which would have satisfied the
law's form while letting anyone through.

\subsection{What the method did and did not find}

Findings 1 to 3 come from the rule and access laws and are about names: two
spellings, two types, two segments treated as one.  Findings 4 to 6 come from
the frontend laws.  Findings 7 to 9 came from asking of every comparison in the
code whether it assumes that the same name means the same bytes, and from
checking each answer against the real file system (including a case-sensitive
volume).  Findings 1, 7 and 8 share one lesson: name equality is a property of
the volume, so deny checks may over-match but allow checks must decide by
identity.  Finding 17 adds a second: a path that the guard has judged is safe
only for code that opens it the way the guard did; another program may see more in
it, so what crosses that boundary must be bytes, not a name.  The laws that held on first check (monotonicity, order
independence, ancestor closure, narrowing, search implies listed, denied is
absent, the round trips) are as informative as the failures, since they turn
long-standing assumptions into checked statements.
